List coloring was introduced independently by Vizing~\cite{Vizing1976}
and by Erd\H{o}s, Rubin, and Taylor~\cite{ErdosRubinTaylor1980}.
Dvo\v{r}\'ak and Postle introduced correspondence coloring in their
2015 preprint, published as~\cite{DvorakPostle}, to allow the
identification of colors to vary from edge to edge. This flexibility
was central to their application to list coloring planar graphs.
An arbitrary correspondence assignment need not admit globally
consistent color names, so a list-coloring argument that uses
independence across such names requires additional work.

An early step toward logarithmic coloring bounds imposed a stronger
restriction on short cycles. Kim~\cite{Kim1995} proved
$\chi_{\mathrm{list}}(G)\le(1+o(1))\Delta/\log\Delta$ for graphs
of girth at least five. Johansson subsequently obtained the list-coloring bounds
$O(\Delta/\log\Delta)$ for triangle-free graphs and
$O_r(\Delta\log\log\Delta/\log\Delta)$ for $K_r$-free graphs.
His triangle-free argument is presented by Molloy and
Reed~\cite[Chapter~13]{MolloyReed2002}; the extended version of
Bansal, Gupta, and Guruganesh~\cite[Sections~6--9]{BansalGuptaGuruganesh2018}
develops his recursive probability modifications and entropy estimates.
Molloy~\cite[Theorems~1 and~2]{Molloy2019} obtained the
triangle-free list bound with leading constant $1+o(1)$ and gave a new
proof of the clique-free bound, with explicit dependence on $r$.

Bernshteyn~\cite{Bernshteyn2016} extended Johansson's triangle-free
bound to correspondence coloring, and later extended Molloy's
asymptotic triangle-free bound and the clique-free bound to that
setting~\cite[Theorems~1.3 and~1.5]{Bernshteyn2019}. In particular,
the previous general upper bound for fixed $r\ge4$ was
\[
 \chi_{\mathrm{DP}}(G)
 =O_r\!\left(\frac{\Delta\log\log\Delta}{\log\Delta}\right).
\]
Davies, Kang, Pirot, and Sereni~\cite[Theorem~30]{DKPSLocalOccupancy}
developed a framework based on local occupancy in the hard-core model
and improved the leading constants. Przyby\l o~\cite[Theorem~10]{Przybylo2024}
also obtained bounds under clique exclusion in the graph whose vertices
are list entries and whose edges are forbidden pairs, while retaining
the maximum degree of the underlying graph as the degree parameter.
These results retain the
$\log\log\Delta$ factor for a general fixed forbidden clique.

Other families of forbidden subgraphs admit the sharper order.
Alon, Krivelevich, and Sudakov~\cite{AKS1999} proved that a graph
whose neighborhoods each span at most $\Delta^2/f$ edges, with
$2\le f\le\Delta^2$, has chromatic number $O(\Delta/\log f)$.
Vu~\cite{Vu2002} extended this local-sparsity bound to list coloring.
Alon, Krivelevich, and Sudakov also deduced
$O_F(\Delta/\log\Delta)$ for every fixed almost bipartite forbidden
graph $F$, meaning that deleting at most one vertex from $F$ leaves a
bipartite graph. Anderson, Bernshteyn, and
Dhawan~\cite{AndersonBernshteynDhawan2023,AndersonBernshteynDhawan}
proved correspondence-coloring bounds with leading constants
$1+o(1)$ for bipartite forbidden graphs and $4+o(1)$ for almost
bipartite forbidden graphs. Their more general theorems assume that
the graph of forbidden pairs excludes $F$ and use its maximum degree
in place of $\Delta(G)$.

More recently, Dhawan, Janzer, and
Methuku~\cite[Theorems~1.4 and~1.5]{DhawanJanzerMethuku} proved
$O_t(\Delta/\log\Delta)$ for ordinary coloring of $K_{t,t,t}$-free
graphs, and hence for every fixed three-colorable forbidden graph.
They also obtained the independence bound
$(1-o(1))n\log d/d$ in terms of the average degree $d$.
Their coloring theorem
concerns ordinary coloring; it does not supply a list or correspondence
bound. Theorem~\ref{thm:main} removes the $\log\log\Delta$ factor
for correspondence coloring with any fixed forbidden clique $K_r$, $r\ge4$.

The structural part of our argument is related to the independence
problem. Ajtai, Erd\H{o}s, Koml\'os, and
Szemer\'edi~\cite{AEKS1981} asked for a lower bound of order
$n\log d/d$ for $K_r$-free graphs of average degree $d$.
Shearer~\cite{Shearer1983,Shearer1995} obtained the asymptotic
triangle-free bound and a clique-free bound with a
$\log\log d$ loss. The companion manuscript
\emph{A logarithmic independence bound for clique-free graphs}
\cite{OpenAIIndependence} proves the logarithmic bound and develops
the weighted multipliers, neighborhood walks, and vertex splitting
used here. The recursive use of clique-free neighborhoods has a
precedent in~\cite[Section~3]{AEKS1981}; entropy increments controlling
distances between walk distributions have a precedent
in~\cite[Section~2, Theorem~8]{BDKY2015}. We adapt these ingredients
to a local adjustment statement that also bounds how many entries
from each list an operation affects. We establish the required
weighted and dual estimates directly.

The remaining probabilistic ingredients also have established
antecedents. We use the simultaneous coupling bound presented by
Angel and Spinka~\cite{AngelSpinka}, with its finite-state proof
included. Concentration for sums in which each independent input affects
only a bounded number of summands is the setting of
Gavinsky, Lovett, Saks, and Srinivasan~\cite{GavinskyReadK}.
We derive the weighted form needed here from Finner's generalized
H\"older inequality~\cite{Finner} and Hoeffding's exponential-moment
bound~\cite{Hoeffding1963}. The Lov\'asz local
lemma~\cite{ErdosLovasz1975} then makes the local round estimates
hold simultaneously.
